\documentclass{beamer}

\usepackage[utf8]{inputenc}
\usepackage{url}
\usepackage{listings}
\usepackage{hyperref}
\usepackage{xcolor}
\usepackage{graphicx}

\input{../helper}


\title{Introdução à Computação \\ uso dos números hexadecimais}
\author{Alexandre Rademaker}
\date{}

\begin{document}

\begin{frame}
 \maketitle
\end{frame}


\begin{frame}[fragile,allowframebreaks]{Uso de hexadecimais}

  Vamos usar hexadecimais para representar endereços da memória.

  \begin{center}
    \includegraphics[width=.7\textwidth]{addresses.png}    
  \end{center}

  \framebreak

  Com 4 bytes (32 bits), só poderíamos contar até cerca de 4 bilhões,
  ou cerca de 4 GB de memória.

  \begin{align*} 
    2^{32} & = 4294967296_{10} \\
          & = 11111111111111111111111111111111_2 \\
          & = FFFFFFFF_{16} 
  \end{align*}
  
  O número hexadecimal é bem mais compacto.

  \framebreak

  Em computadores modernos, usamos 64 bits para endereçar os bytes de
  memória disponíveis (na verdade mais do que podemos efetivamente
  construir):

  \[
    2^{64} = 18.446.744.073.709.551.616 = 16\ exibytes
  \]

  \framebreak

  Os valores na memória de um computador ainda são armazenados como
  binários, mas essa forma de representação nos ajuda a representar
  valores numéricos maiores com menos dígitos.

  Com 8 bits em binário (1 byte), o maior valor que podemos contar é
  255 (11111111). Dois dígitos em hexadecimal podem representar o
  valor de um byte em binário.
  
\end{frame}


\begin{frame}[fragile,allowframebreaks]{Pixels}

  Mas os números hexadecimais também são usados para representar
  cores!

  Aprendemos como arrays representam espaços contínuos da memória
  usados para armazenar valores de um mesmo tipo.

  Também já vimos que toda imagem só pode armazenar um valor finito de
  bits. Cada ponto da imagem chamamos de \textbf{pixel} e uma imagem é
  um array de pixels.

  \framebreak

  Em imagens PB (preto e branco), cada pixel armazena um bit, 1
  (preto) e 0 (branco). 

  \begin{center}
    \includegraphics[width=.7\textwidth]{smiley.png}
  \end{center}

  \framebreak

  Uma imagem pode ser também uma matriz de \ilcode{char} da tabela
  ASCII (caracteres), as chamadas \href{https://www.asciiart.eu}{ASCII
    art}. \vspace{1cm}
  
  \begin{center}
    \begin{lstlisting}[basicstyle=\ttfamily\Large]
     / \__
    (    @\___
    /         O
   /   (_____/
  /_____/   U    
   \end{lstlisting}
  \end{center}

  \framebreak

  Mas também podemos ter imagens coloridas, com uma certa quantidade
  de {\color{red} vermelho}, {\color{green} verde} e {\color{blue}
    azul} por pixel.

  \begin{center}
    \includegraphics[width=.4\textwidth]{eyes.png}
  \end{center}
  
  \framebreak

  Em programas editores de imagens, a escolha de uma certa tonalidade
  de verde, vermelho ou azul está associada a valores em
  hexadecimais. \vspace{.7cm}

  \begin{columns}
    \begin{column}{0.5\textwidth}
      \begin{center}
        \includegraphics[width=.9\textwidth]{color_picker.png}
      \end{center}
    \end{column}
    \begin{column}{0.5\textwidth}
      \begin{center}
        \begin{tabular}{ll} \hline
          code & color \\ \hline
          {\ttfamily \#000000} & black\\
          {\ttfamily \#FFFFFF} & {\color{white} \colorbox{black}{white}}\\
          {\ttfamily \#FF0000} & {\color{red} Red} \\
          {\ttfamily \#00FF00} & {\color{green} Green}\\
          {\ttfamily \#0000FF} & {\color{blue} Blue}\\ \hline
        \end{tabular}
      \end{center}
    \end{column}
  \end{columns}

  Ver \href{https://htmlcolorcodes.com/color-picker/}{seletor de cores} e
  \href{https://en.wikipedia.org/wiki/White}{branco}.

\end{frame}


\end{document}

%%% Local Variables:
%%% mode: latex
%%% TeX-master: t
%%% End:
